\chapter{Communication Complexity — Det Composition}
%%% generated by the blueprint pipeline from TCSlib.CommunicationComplexity.DeterministicCC.DetComposition
\section{Overview}

This module develops the monadic structure of finite-message deterministic protocols:
\texttt{CommunicationComplexity.Deterministic.FiniteMessage.Protocol.map}, \texttt{CommunicationComplexity.Deterministic.FiniteMessage.Protocol.bind}, binary \texttt{CommunicationComplexity.Deterministic.FiniteMessage.Protocol.prod}, and $k$-fold \texttt{CommunicationComplexity.Deterministic.FiniteMessage.Protocol.pi} constructors,
together with their correctness and complexity accounting theorems.

%%%%%%%%%%%%%%%%%%%%%%%%%%%%%%%%%%%%%%%%%%%%%%%%%%%%%%%%%%%%%%%%%%%%%%%%%%%%%%%
\section{Declarations}

\begin{definition}[Map over protocol output]
\lean{CommunicationComplexity.Deterministic.FiniteMessage.Protocol.map}
\label{CommunicationComplexity.Deterministic.FiniteMessage.Protocol.map}
\leanok
\uses{CommunicationComplexity.Deterministic.FiniteMessage.Protocol}
Given a function $g : \alpha \to \beta$ and a protocol $p : \mathrm{Protocol}\,X\,Y\,\alpha$,
\texttt{CommunicationComplexity.Deterministic.FiniteMessage.Protocol.map}~$g$~$p$ is the protocol that runs $p$ and then applies $g$ to the output,
yielding a protocol of type $\mathrm{Protocol}\,X\,Y\,\beta$ with the same communication
structure as $p$.
\end{definition}

\begin{theorem}[Mapping a protocol's output commutes with execution]
\lean{CommunicationComplexity.Deterministic.FiniteMessage.Protocol.map_run}
\label{CommunicationComplexity.Deterministic.FiniteMessage.Protocol.map_run}
\leanok
\uses{CommunicationComplexity.Deterministic.FiniteMessage.Protocol, CommunicationComplexity.Deterministic.FiniteMessage.Protocol.map, CommunicationComplexity.Deterministic.FiniteMessage.Protocol.run}
\difficulty{2}
Let $g : \alpha \to \beta$ be a function between output types, let $p$ be a
finite-message protocol with outputs in $\alpha$, and let $x \in X$ and $y \in Y$ be
inputs. Then \[ \mathrm{run}\bigl(\mathrm{map}(g, p),\, x,\, y\bigr) =
g\bigl(\mathrm{run}(p, x, y)\bigr). \]
\end{theorem}

\begin{theorem}[Post-processing preserves communication complexity]
\lean{CommunicationComplexity.Deterministic.FiniteMessage.Protocol.map_complexity}
\label{CommunicationComplexity.Deterministic.FiniteMessage.Protocol.map_complexity}
\leanok
\uses{CommunicationComplexity.Deterministic.FiniteMessage.Protocol, CommunicationComplexity.Deterministic.FiniteMessage.Protocol.complexity, CommunicationComplexity.Deterministic.FiniteMessage.Protocol.map}
\difficulty{2}
Let $X$, $Y$, $\alpha$, and $\beta$ be types, let $g : \alpha \to \beta$ be any
function, and let $p$ be a finite-message deterministic two-party communication protocol
with input types $X$, $Y$ and output type $\alpha$. Applying $g$ to the output of $p$ —
that is, running $p$ and post-processing its result through $g$ — yields a protocol
whose communication complexity equals that of $p$: \[
\mathrm{complexity}(\mathrm{map}\,g\,p) = \mathrm{complexity}(p). \]
\end{theorem}

\begin{definition}[Monadic bind of protocols]
\lean{CommunicationComplexity.Deterministic.FiniteMessage.Protocol.bind}
\label{CommunicationComplexity.Deterministic.FiniteMessage.Protocol.bind}
\leanok
\uses{CommunicationComplexity.Deterministic.FiniteMessage.Protocol}
Given a protocol $p : \mathrm{Protocol}\,X\,Y\,\alpha$ and a continuation
$q : \alpha \to \mathrm{Protocol}\,X\,Y\,\beta$, \texttt{bind}~$p$~$q$ sequentially
composes them: each output leaf $a$ of $p$ is replaced by the sub-protocol $q\,a$,
giving a protocol of type $\mathrm{Protocol}\,X\,Y\,\beta$.
\end{definition}

\begin{theorem}[Execution of a sequentially composed protocol]
\lean{CommunicationComplexity.Deterministic.FiniteMessage.Protocol.bind_run}
\label{CommunicationComplexity.Deterministic.FiniteMessage.Protocol.bind_run}
\leanok
\uses{CommunicationComplexity.Deterministic.FiniteMessage.Protocol, CommunicationComplexity.Deterministic.FiniteMessage.Protocol.bind, CommunicationComplexity.Deterministic.FiniteMessage.Protocol.run}
\difficulty{2}
Let $p$ be a finite-message protocol with input types $X$, $Y$ and output type $\alpha$,
and let $q$ assign to each value $a \in \alpha$ a finite-message protocol $q(a)$ with
the same input types and output type $\beta$. For any inputs $x \in X$ and $y \in Y$,
executing the sequential composition of $p$ with $q$ agrees with first running $p$ on
$(x,y)$ to obtain an output $a$ and then running $q(a)$ on the same inputs:
\[
\mathrm{run}\bigl(\mathrm{bind}(p,q)\bigr)(x,y) \;=\;
\mathrm{run}\bigl(q(\mathrm{run}(p)(x,y))\bigr)(x,y).
\]
\end{theorem}

\begin{theorem}[Supremum commutes with adding a constant]
\lean{CommunicationComplexity.Deterministic.FiniteMessage.Protocol.finset_sup_add_const}
\label{CommunicationComplexity.Deterministic.FiniteMessage.Protocol.finset_sup_add_const}
\leanok
\difficulty{4}
Let $\iota$ be a finite, nonempty type, let $f : \iota \to \bbn$ be a function, and let
$c \in \bbn$ be a natural number. Then
\[
  \sup_{i \in \iota} \bigl(f(i) + c\bigr) \;=\; \Bigl(\sup_{i \in \iota} f(i)\Bigr) + c.
\]
\end{theorem}

\begin{theorem}[Additivity of protocol complexity under constant-cost bind]
\lean{CommunicationComplexity.Deterministic.FiniteMessage.Protocol.bind_complexity_const}
\label{CommunicationComplexity.Deterministic.FiniteMessage.Protocol.bind_complexity_const}
\leanok
\uses{CommunicationComplexity.Deterministic.FiniteMessage.Protocol, CommunicationComplexity.Deterministic.FiniteMessage.Protocol.bind, CommunicationComplexity.Deterministic.FiniteMessage.Protocol.complexity, CommunicationComplexity.Deterministic.FiniteMessage.Protocol.finset_sup_add_const}
\difficulty{4}
Let $p$ be a finite-message protocol with inputs of types $X$, $Y$ and outputs in a type
$\alpha$, and let $q$ assign to each output value $a \in \alpha$ a finite-message
protocol $q\,a$ with inputs of types $X$, $Y$ and outputs in a type $\beta$. Suppose
there is a natural number $c$ such that every continuation has the same communication
complexity, that is, $\mathrm{complexity}(q\,a) = c$ for all $a \in \alpha$. Then the
sequential composition of $p$ with $q$ has communication complexity
$\mathrm{complexity}(p) + c$.
\end{theorem}

\begin{definition}[Product of two protocols]
\lean{CommunicationComplexity.Deterministic.FiniteMessage.Protocol.prod}
\label{CommunicationComplexity.Deterministic.FiniteMessage.Protocol.prod}
\leanok
\uses{CommunicationComplexity.Deterministic.FiniteMessage.Protocol, CommunicationComplexity.Deterministic.FiniteMessage.Protocol.comap, CommunicationComplexity.Deterministic.FiniteMessage.Protocol.map}
Given $p_1 : \mathrm{Protocol}\,X_1\,Y_1\,\alpha_1$ and
$p_2 : \mathrm{Protocol}\,X_2\,Y_2\,\alpha_2$, their product
$p_1.\texttt{prod}\;p_2 : \mathrm{Protocol}\,(X_1 \times X_2)\,(Y_1 \times Y_2)\,(\alpha_1 \times \alpha_2)$
runs $p_1$ on the first components and $p_2$ on the second components, pairing the outputs.
\end{definition}

\begin{theorem}[Product protocol computes componentwise]
\lean{CommunicationComplexity.Deterministic.FiniteMessage.Protocol.prod_run}
\label{CommunicationComplexity.Deterministic.FiniteMessage.Protocol.prod_run}
\leanok
\uses{CommunicationComplexity.Deterministic.FiniteMessage.Protocol, CommunicationComplexity.Deterministic.FiniteMessage.Protocol.prod, CommunicationComplexity.Deterministic.FiniteMessage.Protocol.run}
\difficulty{2}
Let $p_1$ be a finite-message protocol with input types $X_1$, $Y_1$ and output type
$\alpha_1$, and let $p_2$ be a finite-message protocol with input types $X_2$, $Y_2$ and
output type $\alpha_2$. For every pair of inputs $x = (x_1, x_2) \in X_1 \times X_2$ and
$y = (y_1, y_2) \in Y_1 \times Y_2$, running the product protocol of $p_1$ and $p_2$ on
$(x, y)$ yields the pair of the two componentwise executions:
\[
\mathrm{run}(p_1 \times p_2,\, x,\, y) \;=\; \bigl(\mathrm{run}(p_1,\, x_1,\, y_1),\;
\mathrm{run}(p_2,\, x_2,\, y_2)\bigr).
\]
\end{theorem}

\begin{theorem}[Additivity of complexity under products]
\lean{CommunicationComplexity.Deterministic.FiniteMessage.Protocol.prod_complexity}
\label{CommunicationComplexity.Deterministic.FiniteMessage.Protocol.prod_complexity}
\leanok
\uses{CommunicationComplexity.Deterministic.FiniteMessage.Protocol, CommunicationComplexity.Deterministic.FiniteMessage.Protocol.complexity, CommunicationComplexity.Deterministic.FiniteMessage.Protocol.finset_sup_add_const, CommunicationComplexity.Deterministic.FiniteMessage.Protocol.prod}
\difficulty{4}
Let $p_1$ be a finite-message protocol over inputs $X_1$, $Y_1$ with outputs in
$\alpha_1$, and let $p_2$ be a finite-message protocol over inputs $X_2$, $Y_2$ with
outputs in $\alpha_2$. Then the communication complexity of their product protocol,
which runs $p_1$ on the first input components and $p_2$ on the second and pairs the two
outputs, equals the sum of their individual complexities:
\[
\mathrm{complexity}(p_1 \times p_2) \;=\; \mathrm{complexity}(p_1) +
\mathrm{complexity}(p_2).
\]
\end{theorem}

\begin{definition}[$k$-fold product of protocols]
\lean{CommunicationComplexity.Deterministic.FiniteMessage.Protocol.pi}
\label{CommunicationComplexity.Deterministic.FiniteMessage.Protocol.pi}
\leanok
\uses{CommunicationComplexity.Deterministic.FiniteMessage.Protocol, CommunicationComplexity.Deterministic.FiniteMessage.Protocol.bind, CommunicationComplexity.Deterministic.FiniteMessage.Protocol.comap, CommunicationComplexity.Deterministic.FiniteMessage.Protocol.map}
Given a family of protocols $p_i : \mathrm{Protocol}\,X_i\,Y_i\,\alpha_i$ for
$i \in \mathrm{Fin}\,k$, \texttt{pi}~$p$ is the protocol of type
$\mathrm{Protocol}\,(\Pi_i\,X_i)\,(\Pi_i\,Y_i)\,(\Pi_i\,\alpha_i)$ that runs each $p_i$
on the $i$-th components and collects all outputs into a tuple.
It is defined by induction on $k$, using \texttt{bind}, \texttt{comap}, and \texttt{map}.
\end{definition}

\begin{theorem}[Running the product protocol componentwise]
\lean{CommunicationComplexity.Deterministic.FiniteMessage.Protocol.pi_run}
\label{CommunicationComplexity.Deterministic.FiniteMessage.Protocol.pi_run}
\leanok
\uses{CommunicationComplexity.Deterministic.FiniteMessage.Protocol, CommunicationComplexity.Deterministic.FiniteMessage.Protocol.bind_run, CommunicationComplexity.Deterministic.FiniteMessage.Protocol.comap_run, CommunicationComplexity.Deterministic.FiniteMessage.Protocol.map_run, CommunicationComplexity.Deterministic.FiniteMessage.Protocol.pi, CommunicationComplexity.Deterministic.FiniteMessage.Protocol.run}
\difficulty{4}
Let $k \in \bbn$, and for each index $i \in \mathrm{Fin}\,k$ let $p_i$ be a
finite-message protocol over input types $X_i$, $Y_i$ with output type $\alpha_i$. Given
input tuples $x = (x_i)_{i}$ and $y = (y_i)_{i}$, running the $k$-fold product protocol
on $x$ and $y$ yields the tuple whose $i$-th coordinate is the result of running $p_i$
on the components $x_i$ and $y_i$; that is,
\[
\mathrm{run}\bigl(\Pi_i\, p_i\bigr)\,x\,y \;=\; \bigl(\, i \mapsto
\mathrm{run}(p_i)\,x_i\,y_i \,\bigr).
\]
\end{theorem}

\begin{theorem}[Additivity of communication complexity under products]
\lean{CommunicationComplexity.Deterministic.FiniteMessage.Protocol.pi_complexity}
\label{CommunicationComplexity.Deterministic.FiniteMessage.Protocol.pi_complexity}
\leanok
\uses{CommunicationComplexity.Deterministic.FiniteMessage.Protocol, CommunicationComplexity.Deterministic.FiniteMessage.Protocol.bind_complexity_const, CommunicationComplexity.Deterministic.FiniteMessage.Protocol.comap, CommunicationComplexity.Deterministic.FiniteMessage.Protocol.complexity, CommunicationComplexity.Deterministic.FiniteMessage.Protocol.map, CommunicationComplexity.Deterministic.FiniteMessage.Protocol.pi}
\difficulty{4}
Let $k$ be a natural number, and for each index $i \in \{0, 1, \dots, k-1\}$ let $p_i$
be a finite-message protocol over input types $X_i$, $Y_i$ with output type $\alpha_i$.
The $k$-fold product protocol runs each $p_i$ on the $i$-th coordinates of a pair of
input tuples in $\prod_i X_i$ and $\prod_i Y_i$ and collects the outputs into a tuple in
$\prod_i \alpha_i$. Then the communication complexity of the $k$-fold product protocol
equals the sum of the communication complexities of the factors:
\[
\mathrm{complexity}\!\left(\prod_{i} p_i\right) = \sum_{i=0}^{k-1}
\mathrm{complexity}(p_i).
\]
\end{theorem}
